\documentclass[journal]{IEEEtran}

\ifCLASSINFOpdf
  \usepackage[pdftex]{graphicx}
\else
\fi

\begin{document}
\title{Scalable Data Stream Processor}
\author{Baiqing Lyu \\Boston University}
\maketitle

\begin{abstract}
This document lays out the inspiration for a scalable data stream processor.
It is a general framework for data stream processing, which can be used to
build a wide variety of data processing applications. The uniqueness of this
framework is that it is designed to be scalable, and it is designed to be used
in a distributed environment.
\end{abstract}

\section{Introduction}

\IEEEPARstart{T}{he} current state of the art in data stream processing is software 
frameworks like Apache Flink, Apache Spark, and Apache Storm. These frameworks require
the data pipeline to restart every time a scaling decision has to be made. This inspired
our idea of creating a data stream processor that requires no down time when scaling up
and down, as in a real life production scenario it would happen frequently as high 
and low traffic hours emerge during the day.

\hfill March 31, 2022

\section{Design}

\subsection{Requirements}
Being able to scale up and down without downtime is a fundamental requirement for 
this data stream processor. It would also need to deliver \textbf{at least once} 
processing guarantees when ingesting incoming data. Current scope is to focus on the
feasibility of a data stream processor that can be used in a distributed environment.
\subsection{Goals}
1. Create MVP of a distributed networks that allow nodes to come in and out of the 
network.

2. Give the ability to input data into the network and have a sink that can recieve 
the output.

3. Develop a MVP of a data stream processor that can support easy map reduce operations 
(i.e. wordcount)
\subsection{Software Design}

We will initially use the Chord protocol to manage the way that nodes communicate with each other. 
Chord creates a ring structure for the nodes and provies the ability to map a given key to some node.

As Chord supports the operation of load balancing in addition to mapping keys to nodes and controlling
the behavior when the nodes are down, it is very helpful as the baseline for the data stream processor.
The usage of finger tables within Chord provides a robust way to handle arbritrary node failures. In addition,
this can also handle the scale up and down of nodes within the processor.

The main purpose of Chord is very simple in its mission, \textit{provided a key, the protocol will
return the node that is responsible for the key}. We use this property to implement the map reduce functionality
and route events to their correct nodes.

\subsection{Scaling}

\subsubsection*{Phase 1}
\subsubsection*{Phase 2}
\subsubsection*{Phase 3}

\end{document}